\begin{table}[htbp]
\centering
\caption{Parámetros físicos, geométricos y electromagnéticos nominales empleados en la simulación computacional del experimento de Rimstar.}
\label{tab:parametros}
\begin{ruledtabular}
\begin{tabular}{lccc}
\textbf{Parámetro} & \textbf{Símbolo} & \textbf{Valor Nominal} & \textbf{Incertidumbre} \\
\hline
Masa del imán & $m_0$ & $28.50\ \mathrm{g}$ & $\pm 0.20\ \mathrm{g}$ \\
Radio del imán & $a$ & $7.50\ \mathrm{mm}$ & $\pm 0.05\ \mathrm{mm}$ \\
Longitud del imán & $2h$ & $10.00\ \mathrm{mm}$ & $\pm 0.05\ \mathrm{mm}$ \\
Momento dipolar efectivo & $m_{\mathrm{eff}}$ & $0.820\ \mathrm{A\cdot m^2}$ & $\pm 0.025\ \mathrm{A\cdot m^2}$ \\
Longitud del tubo & $L$ & $1.000\ \mathrm{m}$ & $\pm 0.002\ \mathrm{m}$ \\
Radio medio del tubo & $R$ & $11.50\ \mathrm{mm}$ & $\pm 0.10\ \mathrm{mm}$ \\
Espesor de pared & $\Delta R$ & $1.20\ \mathrm{mm}$ & $\pm 0.05\ \mathrm{mm}$ \\
Conductividad Cobre & $\sigma_{\mathrm{Cu}}$ & $5.96 \times 10^7\ \mathrm{S/m}$ & $\pm 0.06 \times 10^7\ \mathrm{S/m}$ \\
Conductividad Aluminio & $\sigma_{\mathrm{Al}}$ & $3.77 \times 10^7\ \mathrm{S/m}$ & $\pm 0.05 \times 10^7\ \mathrm{S/m}$ \\
Conductividad PVC (control) & $\sigma_{\mathrm{PVC}}$ & $\approx 10^{-12}\ \mathrm{S/m}$ & --- \\
Aceleración gravitacional & $g$ & $9.80665\ \mathrm{m/s^2}$ & --- \\
\end{tabular}
\end{ruledtabular}
\end{table}
